\documentclass[10pt]{article}

\usepackage[utf8]{inputenc}

\usepackage[T1]{fontenc}

\usepackage{amsmath}

\usepackage{amsfonts}

\usepackage{amssymb}

\usepackage[version=4]{mhchem}

\usepackage{stmaryrd}

\usepackage{mathrsfs}



\begin{document}

Лагранжиан самого поля Янга-Миллса имеет вид:



\[

\mathscr{L}=\frac{1}{8} \operatorname{tr} F_{\mu \nu} F_{\mu \nu}=\frac{1}{8}\left\langle F_{\mu v}, F_{\mu \nu}\right\rangle,

\]



где через $\langle,>$ обозначена форма Киллинга для алгебры Ли группы $\Omega, F_{\mu \nu}-$ тензор напряженности поля Янга-Миллса:



\[

F_{\mu \nu}=\partial_{\mu} A_{v}-\partial_{\nu} A_{\mu}+g\left[A_{\mu}, A_{v}\right]

\]



Если ограничиться минимальным числом производных, то форма (9) является единственным лоренц-инвариантным выражением, инвариантным относительно неабелевых калибровочных преобразований (6).



Уравнения Эйлера - Лагранжа для поля Янга - Миллса имеют вид:



\[

\nabla_{\mu} F_{\mu \nu}=j_{v}

\]



где $j_{V}$ - ток полей материи.



Естественный геометрич. язык для описания полей Янга-Миллса дает теория расслоенных пространств. Полю Янга-Миллса в этой теории соответствует связность в главном расслоении. Подобно символам Кристоффеля в теории тяготения, К. П. описывает параллельный перенос в зарядовом пространстве. Если $\gamma(s)$ - контур в пространстве-времени, задаваемый уравнениями $x_{\mu}=x_{\mu}(s)$, то ковариантная производная в касательном направлении к этому контуру равна нулю:



\[

\left.\nabla_{\mu} \Psi(x)\right|_{x=x(s)} \frac{d x_{\mu}}{d s}=0 .

\]



Тензор напряженности $F_{\mu \nu}$ является тензором кривизны зарядового пространства.



Инвариантность функции Лагранжа относительно калибровочных преобразований, зависящих от произвольной функции, приводит к линейной зависимости уравнений движения. Нек-рые из функций, параметризующих классич. решение, произвольным образом зависят от времени. При квантовании необходимо разделить истинные динамич. переменные и групповые параметры. Так, для поля Янга-Миллса компонента $a$ представляет собой не динамич. переменную, а множитель Лагранжа. Соответствующий ей канонич. импульс



\[

p_{0}^{a}=\delta \varrho / \delta\left(\frac{\partial A^{a}}{\partial x_{0}}\right)

\]



тождественно обращается в нуль, а уравнение Эйлера - Лагранжа, получающееся при варьировании действия по $A_{0}$ :



\[

\nabla_{i} p_{i}=j_{0}, i=1,2,3,

\]



не содержит производной по времени и поэтому не описывает динамику системы, а является уравнением связи. Общая теория квантования систем со связями была развита П. Дираком (P. Dirac, 1956), Л. Д. Фаддеевым (1969) и др. Ее суть состоит в том, что канонич. перестановочные соотношения накладываются лишь на истинные динамич. переменные, к-рые можно найти, решив уравнения связи и наложив на поля $A_{\mu}$ дополнительные условия, выбирающие единственного представителя в классе калибровочно эквивалентных конфигураций $A_{\mu}^{\omega}$. Эти условия называются у с ло в и я м и ка л и б ро в к и (см. Калибровочное условие). Если, напр., наложить на поля $A_{\mu}$ условие кулоновской калибровки $\partial_{i} A_{i}=0$, то уравнения связи (13) можно явно решить, выразив продольную часть вектора импульса $p_{i}$ через трехмерно поперечные компоненты $p_{i}^{t r}, A_{i}^{t r}$. Если подставить это решение в исходное действие, то оно будет зависеть лишь от поперечных компонент $p_{i}^{t r}, A_{i}^{t r}$, к-рые и являются в данном случае истинными динамич. переменными. На них можно наложить канонич. перестановочные соотношения. В электродинамике подобная процедура соответствует описанию системы в терминах поперечно поляризованных фотонов. В случае неабелевых К. п. решение уравнения (13) представляет собой бесконечный ряд по константе связи $g$, подстановка к-рого в действие порождает нелинейные по полям члены, отвечающие нулевым вершинам взаимодействия. Релятивистски инвариантная теория возмушений для квантованного поля Янга - Миллса была развита Л. Д. Фаддеевым и В. Н. Поповым (1967) и Б. Де-Виттом (B. De-Witt, 1967) с помощью метода континуального интеграла. Фейнмановская диаграммная техника в этом случае может быть описана эффективным действием вида



\[

\begin{aligned}

S_{\ni \Phi} & =\int\left\{\frac { 1 } { 8 } \operatorname { T r } \left[F_{\mu \nu} F_{\mu \nu}-\frac{2}{\beta}\left(\partial_{\mu} A_{\mu}\right)^{2}-\right.\right. \\

& \left.\left.-4 \bar{c}\left(\square c-g \partial_{\mu}\left[A_{\mu}, c\right]\right)\right]\right\} d^{4} x .

\end{aligned}

\]



Здесь $\beta$ - параметр, фиксирующий калибровку, $\bar{c}, c-$ скалярные поля, подчиняющиеся статистике Ферми - Дирака. Эти поля называются духами Фаддеева- Попова. Они участвуют лишь во внутренних линиях фейнмановских диаграмм и отсутствуют в наблюдаемых асимптотич. соотношениях. Эффективное действие (14) содержит также нефизич. компоненты поля $A_{\mu}$, отвечающие продольным и временным квантам. Роль духовых полей состоит в том, чтобы скомпенсировать вклад этих квантов и обеспечить унитарность матрицы рассеяния в пространстве физич. асимптотич. состояний.



Действие (14) инвариантно относительно преобразований, затрагивающих помимо К. п. $A_{\mu}$ духовые поля $c, \bar{c}$. Эта симметрия, заменяющая калибровочную инвариантность, получила название Б Р С-с им метр и и (С. Becchi, A. Ronet, R. Stora, 1976). БРС-симметрия лежит в основе альтернативного метода квантования К. п., при к-ром теория с самого начала рассматривается в расширенном пространстве, содержащем помимо физич. квантов духовые поля, а также продольные и временные кванты, а векторы физич. подпространства $|\Phi\rangle_{\text {физ выделяются условием }}$



\[

Q|\Phi\rangle_{\text {физ }}=0 \text {, }

\]



где $Q$ - генератор БРС-преобразований.



Для вычисления вероятностей различных процессов необходимо провести процедуру перенормировки функций Грина К. п., устраняющую расходимости за счет перераспределения масс, зарядов и постоянных нормировки волновых функций. При этом наблюдаемые величины не должны зависеть от конкретного выбора условия калибровки и перенормированная матрица рассеяния должна быть унитарна в пространстве физич. состояний. Необходимым и достаточным условием согласованности перенормировочной процедуры с калибровочной инвариантностью является выполнение соотношений между функциями Грина, к-рые получили название тожде ст в Сла внова-Те йл ора (А. А. Славнов, 1972, J. С. Taylor, 1972). Эти соотношения обобщают Уорда тождества в электродинамике на случай неабелевых калибровочных полей. Они связаны с БРС-симметрией эффективного действия.



К. П. лежат в основе большинства современных моделей взаимодействия элементарных частиц. Их успешное применение в физике сильных взаимодействий обусловлено в первую очередь явлением асимптотич. свободы. При нек-рых ограничениях на калибровочную группу и число мультиплетов полей материи эффективное взаимодействие полей Янга - Миллса убывает на малых расстояниях или, что то же самое, при больших энергиях. Это явление, получившее название ас и мпт о т ичес кой с вободы, подтверждается целым рядом экспериментов, в частности экспериментами по глубоконеупругому рассеянию.



В то же время на больших расстояниях эффективное взаимодействие неабелевых К. п. растет и теория возмущений по константе связи непримеініма. Формально это проявляется в том, что интегралы, соответствуюшие диаграммам Фейнмана, расходятся при малых импульсах (см. Инфракрасные расходимости). В отличие от электродинамики, где сушествует регулярный метод устранения этих расходимостей, в неабелевых калибровочных теориях подобная процедура отсутствует. Для исследования калибровочных тео-





\end{document}